% Olympus — Introduction (USENIX ATC-style)
% Drop-in section. Citation keys are placeholders; wire them to your .bib.
\section{Introduction}

Congestion control (CC) decides how fast each sender injects traffic into a
shared network, and in doing so determines the throughput, latency, and
fairness that applications ultimately experience. For three decades the
dominant approach has been hand-crafted heuristics: Reno, CUBIC~\cite{cubic},
and more recently BBR~\cite{bbr} that map a fixed set of congestion signals
onto a fixed reaction. Designing a heuristic that behaves well across the
diversity of modern paths, from data-center fabrics to cellular, satellite, and
wide-area links, is notoriously hard, and each new network regime tends to
expose a new failure mode.

Reinforcement learning offers an appealing alternative. RL-based congestion
control (RL-CC) casts the problem as sequential decision making: a policy
repeatedly observes the state of a flow and adjusts its sending rate or
congestion window to maximize a reward that trades off throughput, delay, and
fairness. A growing body of work---Aurora~\cite{aurora}, Orca~\cite{orca},
Sage~\cite{sage}, Astraea~\cite{astraea}, and others---has shown that learned
policies can match or surpass hand-tuned schemes, and that they can adapt to
conditions their designers never explicitly anticipated.

Yet progress in RL-CC is far harder than this steady cadence of results
suggests. Unfortunately, each implementation comes with its own purpose-built
framework, meaning that significant effort has to go into either altering an
existing framework to support a new problem formulation---such as different
environments, actions, or states---or building a new framework altogether. An
RL-CC framework bakes in a large set of entangled design decisions: which
algorithm trains the policy, how network state is represented, what the action
space looks like, how reward is computed, how flows are emulated, and how
traffic is generated. Existing frameworks fix all of these at design time and
weld them into a single artifact tuned to reproduce one paper's results, rather
than a substrate for the next one. Asking a different research question---a new
state representation, a fairness-aware reward, a more sample-efficient
learner---means either invasive surgery on someone else's codebase or starting
from scratch. The barrier to entry is high, cross-scheme comparisons are
confounded by incompatible substrates, and good ideas are slow to compose.

Three consequences of this status quo are especially limiting. First, existing
frameworks are single-agent. They train one policy controlling one flow in
isolation and treat every other flow as exogenous background traffic. This
leaves no clean path to study the multi-agent settings---where flows jointly
learn to share a bottleneck---that are central to the very properties
(fairness, convergence, responsiveness under churn) congestion control exists
to provide. Second, they are pinned to the learning algorithm that was current
when they were written. Because the learner is fused into the framework,
advances that have reshaped RL---model-based agents and world models,
multi-agent transformers, more sample-efficient off-policy methods---cannot be
adopted without effectively rebuilding the framework around them. Third,
inference is tied to the specific traffic generator and emulator used during
training, so a trained policy is difficult to take anywhere beyond the harness
that produced it.

 We present Olympus, a unified and modular framework for
RL-CC research and deployment. Olympus treats the state, action, reward,
environment, and learner as five independent, swappable components selected per
training run. Changing any one of them---a different observation, a different
congestion-window action mapping, a fairness-oriented reward, a new emulation
backend, or an entirely different learning algorithm---is a configuration
change rather than a fork, and adding a new one is a self-contained extension
rather than a framework rewrite. A single long-lived learner is fed by many
parallel flow rollouts, so the same design serves a single agent and a
population of jointly trained agents without special-casing the orchestration.
The result is a framework on which a new problem formulation can be expressed
and run in minutes, lowering the barrier to entry and enabling the rapid
iteration that purpose-built frameworks foreclose.

A second goal of Olympus is to make trained RL-CC policies easier to use beyond
the specific traffic generators and experimental setups in which they were
developed. Rather than tying inference to a purpose-built traffic-generation
application or a single emulation backend, Olympus separates the policy
execution path from both the workload and the environment used for training or
evaluation. This allows the same trained policy to be evaluated across
different classes of environments, from standard tools such as
Mahimahi~\cite{mahimahi} or Mininet~\cite{mininet}, to discrete-event network
simulators such as OMNeT++~\cite{omnetpp}, to alternative emulators, testbeds,
or real network deployments, while also supporting different traffic generators
and real applications. As a result, Olympus is better suited for both
controlled experimentation and practical deployment studies.

Concretely, Olympus supports the following capabilities, which existing
state-of-the-art RL-CC training frameworks do not provide together:

\begin{itemize}
  \item Single-agent and multi-agent training. To our knowledge, Olympus is the
        first RL-CC framework to support both a single learned flow and multiple
        flows that jointly learn to share a bottleneck, under one orchestration
        model.
  \item State-of-the-art, hot-swappable learners. Olympus ships modern learners
        as interchangeable plugins---model-based agents and world models
        (DreamerV3~\cite{dreamerv3}, and a multi-agent world model), multi-agent
        transformers (MAT~\cite{mat}), off-policy actor-critics
        (TD3~\cite{td3}), model-based policy optimization (MBPO~\cite{mbpo}),
        recurrent PPO, and option-critic---rather than a single algorithm fixed
        at design time.
  \item Swappable state representations. Raw and filtered observations (e.g.\
        Kalman-filtered, min-RTT, oracle, and Orca-style states) are selected
        per run and shared consistently between training and inference.
  \item Swappable action spaces. Congestion-window action mappings (e.g.\ an
        exponential multiplier and Astraea-style relative updates) are chosen
        independently of the learner.
  \item Swappable reward functions. Throughput--delay and fairness-oriented
        objectives, including team and symmetric fair-share rewards for the
        multi-agent setting, are plugins rather than hard-coded terms.
  \item Swappable environments and emulation backends. Emulators and
        simulators---Mininet today, with Mahimahi, OMNeT++, testbeds, and real
        deployments behind the same interface---are selected per run.
  \item A decoupled policy-execution path. The same checkpoint runs in training,
        benchmarking, and deployment without binding to a bespoke traffic
        generator or a single backend.
  \item Parallel rollouts with a shared learner and built-in, config-driven
        benchmark suites for reproducible, apples-to-apples comparison across
        schemes.
\end{itemize}

In summary, this paper makes the following contributions. (i)~We identify the
fragmentation of RL-CC research into incompatible, purpose-built frameworks as
a first-class obstacle to progress, and we characterize the design decisions
that such frameworks needlessly entangle. (ii)~We present Olympus, a modular
framework in which states, actions, rewards, environments, and learners are
independent components that can be changed from one training run to the next and
extended in isolation, and which is, to our knowledge, the first to support both
single-agent and multi-agent RL-CC with state-of-the-art learners. (iii)~We
show how decoupling the policy-execution path from the workload and the
environment lets a single trained policy move across emulators, simulators,
testbeds, and real deployments, narrowing the gap between RL-CC research and
practice.
